\pgfmathparse{0.2ex*1em/1ex}
\tikzset{pxislot/.style={line width=\pgfmathresult, baseline=0.34em/2}}

\newcommand{\pxislotone}{\noindent
  \beginpgfgraphicnamed{Graphics-pxislot1}
  \begin{tikzpicture}[pxislot]
    \draw (0,0) -- (0.5em,0.866em) -- (1em,0) -- (0,0) -- cycle;
    \pgfusepath{use as bounding box}
    \node[scale=0.8] at (0.5em,0.34em) {\texttt{1}};
  \end{tikzpicture}
  \endpgfgraphicnamed
}

\newcommand{\pxislottwo}{\noindent
  \beginpgfgraphicnamed{Graphics-pxislot2}
  \begin{tikzpicture}[pxislot]
    \draw (0.5em,0.5em) circle (0.355em);
    \draw (0,0.5em) -- (0.5em,1em) -- (1em,0.5em) -- (0.5em,0) -- cycle;
    \node[scale=0.8] at (0.5em,0.5em) {\texttt{2}};
  \end{tikzpicture}
  \endpgfgraphicnamed
}

\newcommand{\pxislotn}[1]{\noindent
  \beginpgfgraphicnamed{Graphics-pxislot#1}
  \begin{tikzpicture}[pxislot]
    \draw (0.5em,0.5em) circle (0.5em);
    \node[scale=0.8] at (0.5em,0.5em) {\texttt{#1}};
  \end{tikzpicture}
  \endpgfgraphicnamed
}

% \pxislotone
% \pxislottwo
% \foreach \x in {3,...,7}
% {
%   \pxislotn{\x}
% }


%%% Local Variables: 
%%% mode: latex
%%% TeX-master: "Manual"
%%% End: 
